% ===============================================================
% TEMPLATE DE APRESENTAÇÃO (BEAMER) - TCC1 - SENAC SANTO AMARO
% ===============================================================
% Apresentação da PROPOSTA de Trabalho de Conclusão de Curso (TCC1)
%
% >>> DURAÇÃO DA APRESENTAÇÃO: 20 MINUTOS <<<
%
% A apresentação DEVE seguir a mesma estrutura do documento escrito
% (arquivo tcc.tex), na ordem:
%   1. Introdução (Contexto, Justificativa, Objetivos)
%   2. Referencial Bibliográfico (Metodologia da pesquisa,
%      Referencial, Trabalhos relacionados)
%   3. Desenvolvimento (Solução proposta, Cronograma)
%
% LOGO: a logo do SENAC aparece GRANDE no primeiro slide (capa) e
% PEQUENA no canto superior direito dos demais slides.
% Coloque o arquivo da logo em:  assets/logo-senac.png
% (Se o arquivo não existir, um marcador de texto é exibido no lugar
%  para que a apresentação compile mesmo assim.)
%
% COMO COMPILAR:
%   pdflatex apresentacao_tcc1.tex
%   pdflatex apresentacao_tcc1.tex   (2x para fixar índice/numeração)
% ===============================================================

\documentclass[aspectratio=169,12pt,brazil]{beamer}
\usepackage[T1]{fontenc}
\usepackage[utf8]{inputenc}
\usepackage[brazil]{babel}
\usepackage{graphicx}
\usepackage{booktabs}
\usepackage{array}
\usepackage{tikz}
\usetikzlibrary{positioning,arrows.meta}

\usetheme{default}
\useinnertheme{circles}
\setbeamertemplate{navigation symbols}{}
\definecolor{senacblue}{RGB}{0,51,102}
\definecolor{senacorange}{RGB}{243,146,0}
\setbeamercolor{structure}{fg=senacblue}
\setbeamercolor{frametitle}{fg=white,bg=senacblue}
\setbeamercolor{title}{fg=senacblue}
\setbeamercolor{block title}{fg=white,bg=senacblue}
\setbeamercolor{block body}{bg=senacblue!6}
\setbeamerfont{frametitle}{series=\bfseries,size=\normalsize}
\setbeamertemplate{footline}{%
  \leavevmode\hbox{%
    \begin{beamercolorbox}[wd=.32\paperwidth,ht=2.6ex,dp=1ex,leftskip=1em]{author in head/foot}%
      \insertshortauthor
    \end{beamercolorbox}%
    \begin{beamercolorbox}[wd=.52\paperwidth,ht=2.6ex,dp=1ex,center]{title in head/foot}%
      \insertshorttitle
    \end{beamercolorbox}%
    \begin{beamercolorbox}[wd=.16\paperwidth,ht=2.6ex,dp=1ex,rightskip=1em]{title in head/foot}%
      \hfill\insertframenumber/\inserttotalframenumber
    \end{beamercolorbox}%
  }%
}
\setbeamercolor{author in head/foot}{bg=senacblue,fg=white}
\setbeamercolor{title in head/foot}{bg=senacblue!85,fg=white}

\newcommand{\logosenacgrande}{%
  \IfFileExists{assets/logo-senac.png}{%
    \includegraphics[height=1.8cm]{assets/logo-senac.png}%
  }{\fcolorbox{senacblue}{white}{\parbox[c][1.2cm][c]{3.3cm}{\centering\Large\bfseries\color{senacblue}SENAC}}}%
}
\newcommand{\logosenacpequena}{%
  \IfFileExists{assets/logo-senac.png}{%
    \includegraphics[height=.45cm]{assets/logo-senac.png}%
  }{\fcolorbox{senacblue}{white}{\scriptsize\bfseries\color{senacblue}SENAC}}%
}
% Caixas de largura fixa e barras para os dois diagramas da apresentação.
\newcommand{\caixafluxo}[1]{%
  \fcolorbox{senacblue}{senacblue!7}{%
    \parbox[c][1.15cm][c]{2.85cm}{\centering\scriptsize #1}}%
}
\newcommand{\barrames}{\textcolor{senacblue}{\rule{1.10cm}{.15cm}}}
\newcommand{\marcocronograma}{\textcolor{senacorange}{$\bullet$}}
\addtobeamertemplate{frametitle}{}{%
  \begin{tikzpicture}[remember picture,overlay]
    \node[anchor=north east,fill=white,rounded corners=2pt,inner sep=2pt]
      at ([xshift=-5pt,yshift=-5pt]current page.north east)
      {\logosenacpequena};
  \end{tikzpicture}%
}

\title[Classificação de FA em ECG]{Comparação entre CNN-BiLSTM e BiLSTM na classificação de fibrilação atrial em sinais de eletrocardiograma}
\subtitle{Análise do impacto da Cross-Entropy e da Focal Loss\\[.2em]Proposta de Trabalho de Conclusão de Curso — TCC1}
\author[Daniel e Eduardo]{Daniel de Araujo Carvalho\\Eduardo Dias Santos}
\institute[SENAC Santo Amaro]{SENAC — Unidade Santo Amaro\\Bacharelado em Ciência da Computação\\Orientador: Prof. Afonso Lelis}
\date{São Paulo, 2026}

\begin{document}

{\setbeamertemplate{footline}{}
\begin{frame}[plain]
  \centering
  \logosenacgrande\par\vspace{.35cm}
  {\large\bfseries\color{senacblue}\inserttitle\par}\vspace{.18cm}
  {\small\insertsubtitle\par}\vspace{.4cm}
  {\small\insertauthor\par}\vspace{.2cm}
  {\footnotesize\insertinstitute\par}\vfill
  {\footnotesize\insertdate\par}
\end{frame}}

\begin{frame}{Roteiro}
  \tableofcontents
\end{frame}

\section{Introdução}

\begin{frame}{Contexto: ECG e fibrilação atrial}
  \begin{itemize}
    \item O eletrocardiograma (ECG) registra a atividade elétrica cardíaca.
    \item A fibrilação atrial (FA) é um ritmo irregular que pode ser identificado no ECG; a análise de registros curtos motiva a classificação automática.
    \item O desafio PhysioNet/CinC 2017 fornece ECG de derivação única: 8.528 registros de treinamento, entre eles 5.154 normais e 771 com FA.
  \end{itemize}
  \begin{block}{Recorte do TCC}
    Classificação binária de 5.925 registros Normal/FA; aproximadamente 6,7 registros normais para cada registro com FA.
  \end{block}
  {\scriptsize Fonte: Clifford et al. (2017).}
\end{frame}

\begin{frame}{Pergunta e hipótese de pesquisa}
  \begin{block}{Pergunta}
    Em que medida a Focal Loss, em comparação à Cross-Entropy, influencia a classificação de FA em ECG nas arquiteturas CNN-BiLSTM e BiLSTM?
  \end{block}
  \begin{block}{Hipótese a testar}
    A Focal Loss melhora o F1 da classe FA em mais de \textbf{2 pontos percentuais} em relação à Cross-Entropy \textbf{em cada arquitetura}.
  \end{block}
  \small A comparação será pareada e adotará nível de significância de 5\%. O limiar de 2 pontos percentuais é um critério de relevância definido para este estudo.
\end{frame}

\begin{frame}{Justificativa e diferencial}
  \begin{itemize}
    \item Trabalhos recentes combinam convolução, recorrência, técnicas para desbalanceamento e funções de perda em \textbf{classificação de batimentos}.
    \item Falta verificar, sob um mesmo protocolo, o efeito da função de perda sobre o \textbf{ritmo FA de um registro curto}.
    \item A comparação CNN-BiLSTM × BiLSTM examina a contribuição da etapa convolucional com processamento bidirecional em ambos os modelos.
    \item Tempo de treinamento, latência e parâmetros serão analisados junto ao desempenho de classificação.
  \end{itemize}
  {\scriptsize Fontes: Al-Mousa et al. (2023); Efe e Yavsan (2024); Albarrak, Alharbi e Ibrahim (2025); Clifford et al. (2017).}
\end{frame}

\begin{frame}{Objetivos}
  \begin{block}{Objetivo principal}
    Comparar experimentalmente CNN-BiLSTM e BiLSTM na classificação de FA, avaliando o impacto da Cross-Entropy e da Focal Loss no desempenho e no custo computacional.
  \end{block}
  \begin{enumerate}
    \item Preparar os registros Normal/FA do PhysioNet/CinC 2017.
    \item Implementar e treinar as quatro combinações.
    \item Medir F1 de FA e demais métricas de classificação.
    \item Comparar o efeito das funções de perda e o custo das arquiteturas.
  \end{enumerate}
\end{frame}

\section{Referencial Bibliográfico}

\begin{frame}{Metodologia da pesquisa bibliográfica}
  \begin{itemize}
    \item Levantamento exploratório em registros do PubMed, documentação e artigo do PhysioNet/CinC 2017 e periódicos da área.
    \item Termos de busca: \texttt{"atrial fibrillation" AND "single-lead ECG" AND "deep learning"}; \texttt{"focal loss" AND "ECG" AND "class imbalance"}.
    \item Artigos relacionados originais, em inglês, publicados entre 2023 e 2025; fontes clássicas para os conceitos fundamentais.
    \item Separação entre trabalhos sobre \textbf{batimentos} e a tarefa deste estudo, sobre \textbf{ritmo por registro}.
  \end{itemize}
\end{frame}

\begin{frame}{Fundamentos: do sinal às arquiteturas}
  \begin{itemize}
    \item \textbf{ECG e FA:} irregularidade dos intervalos RR e ausência de ondas P discerníveis são características relevantes.
    \item \textbf{CNN 1D:} extrai padrões locais diretamente da sequência de amostras.
    \item \textbf{BiLSTM:} processa dependências temporais nos dois sentidos da sequência.
    \item \textbf{CNN-BiLSTM:} acrescenta extração convolucional antes da mesma abordagem recorrente bidirecional.
  \end{itemize}
  {\scriptsize Referencial do texto: Clifford et al. (2017); Goodfellow, Bengio e Courville (2016); Hochreiter e Schmidhuber (1997); Schuster e Paliwal (1997).}
\end{frame}

\begin{frame}{Funções de perda e classe minoritária}
  \begin{columns}[T]
    \column{.49\textwidth}
    \begin{block}{Cross-Entropy}
      \centering $L_{\mathrm{CE}}(p_t)=-\log(p_t)$
    \end{block}
    \small Referência para o treinamento das duas arquiteturas.
    \column{.49\textwidth}
    \begin{block}{Focal Loss}
      \centering $L_{\mathrm{FL}}(p_t)=-\alpha_t(1-p_t)^\gamma\log(p_t)$
    \end{block}
    \small Reduz a contribuição relativa dos exemplos fáceis; o benefício para FA será testado.
  \end{columns}
  \vspace{.25cm}
  \small $p_t$ é a probabilidade atribuída à classe verdadeira; $\gamma$ controla a ênfase nos exemplos difíceis e $\alpha_t$ o peso da classe.
  \vfill {\scriptsize Fonte: Lin et al. (2017).}
\end{frame}

\begin{frame}{Como medir o desempenho}
  \begin{block}{Métrica principal: F1 da classe FA}
    \centering $F1=\dfrac{2\,VP}{2\,VP+FP+FN}$\\
    \scriptsize $VP$: verdadeiros positivos; $FP$: falsos positivos; $FN$: falsos negativos.
  \end{block}
  \begin{itemize}
    \item \textbf{Revocação:} proporção dos registros com FA detectados.
    \item \textbf{Precisão:} proporção de previsões de FA corretas.
    \item Também: matriz de confusão, acurácia, especificidade, AUC-ROC e AUC-PR.
  \end{itemize}
  \small A acurácia isolada pode esconder dificuldades com a classe menos frequente.
\end{frame}

\begin{frame}{Trabalhos relacionados: tarefas distintas}
  \scriptsize
  \begin{tabular}{p{2.1cm}p{3.05cm}p{3.7cm}p{3.1cm}}
    \toprule
    \textbf{Estudo} & \textbf{Tarefa/base} & \textbf{Método} & \textbf{Diferença para este TCC} \\
    \midrule
    Al-Mousa et al. (2023) & 5 classes de batimentos; MIT-BIH supraventricular & Métodos tradicionais e SMOTE & Não classifica FA por registro \\
    Efe e Yavsan (2024) & 5 classes de batimentos; MIT-BIH Arrhythmia & CNN-BiLSTM, atenção e Focal Loss & Não avalia FA por registro \\
    Albarrak et al. (2025) & 5 classes de batimentos; MIT-BIH & CNN-LSTM e otimização GWO & Não compara o fatorial proposto \\
    \midrule
    \textbf{Este TCC} & \textbf{Normal × FA; CinC 2017} & \textbf{CNN-BiLSTM/BiLSTM × CE/FL} & \textbf{F1 de FA e custo, mesmo protocolo} \\
    \bottomrule
  \end{tabular}
  \vspace{.2cm}

  \small Os resultados numéricos dos estudos sobre batimentos não são diretamente comparáveis aos de FA por registro.
\end{frame}

\section{Desenvolvimento}

\begin{frame}{Solução proposta: fluxo experimental}
  \centering
  \resizebox{.93\linewidth}{!}{%
    \begin{tabular}{@{}ccccc@{}}
      \caixafluxo{CinC 2017\\ECG, 300 Hz} & $\longrightarrow$ &
      \caixafluxo{Seleção Normal/FA\\5.925 registros} & $\longrightarrow$ &
      \caixafluxo{Divisão estratificada\\70\% / 15\% / 15\%} \\[.18cm]
      & & & & $\downarrow$ \\[.18cm]
      \caixafluxo{Avaliação\\F1, estatística e custo} & $\longleftarrow$ &
      \caixafluxo{Treinamento\\quatro configurações} & $\longleftarrow$ &
      \caixafluxo{Pré-processamento\\30 s; Z-score} \\
    \end{tabular}%
  }
  \vspace{.42cm}

  \small As classes “outro ritmo” e “ruidoso” não integram o recorte binário.
\end{frame}

\begin{frame}{Dados, divisão e pré-processamento}
  \begin{itemize}
    \item \textbf{Uma derivação, 300 Hz:} cada entrada terá 30 segundos, ou \textbf{9.000 amostras}.
    \item Registros longos: recorte central. Registros curtos: preenchimento com zeros.
    \item Divisão por registro, estratificada: \textbf{70\% treino, 15\% validação, 15\% teste}; idêntica para as quatro configurações.
    \item Z-score: média e desvio-padrão calculados somente com amostras válidas do treino; aplicar aos demais conjuntos e preencher com zeros \textbf{depois}.
  \end{itemize}
  \small Ferramentas previstas: Python, WFDB, PyTorch, NumPy, Pandas, scikit-learn, SciPy e Matplotlib.
\end{frame}

\begin{frame}{Arquiteturas e funções de perda}
  \begin{columns}[T,onlytextwidth]
    \column{.49\textwidth}
    \begin{block}{BiLSTM}
      \footnotesize
      \begin{itemize}
        \item ECG (9.000 amostras) $\rightarrow$ Average Pooling (fator 4).
        \item Duas camadas BiLSTM: 64 unidades por direção; dropout de 0,2 entre elas.
        \item Camada linear com uma saída.
      \end{itemize}
      \centering \textbf{133.761 parâmetros previstos}
    \end{block}
    \column{.49\textwidth}
    \begin{block}{CNN-BiLSTM}
      \footnotesize
      \begin{itemize}
        \item Conv1D: 16 filtros, kernel 7, padding 3 e ReLU; Average Pooling (fator 4).
        \item As mesmas duas camadas BiLSTM: 64 unidades por direção; dropout de 0,2.
        \item Camada linear com uma saída.
      \end{itemize}
      \centering \textbf{141.569 parâmetros previstos}
    \end{block}
  \end{columns}
  \vspace{.15cm}
  \small \textbf{Em ambas:} CE sem pesos ou Focal Loss com $\gamma=2$, $\alpha_{\mathrm{FA}}=0{,}75$ e $\alpha_{\mathrm{Normal}}=0{,}25$.
\end{frame}

\begin{frame}{Experimento fatorial 2 × 2}
  \centering
  \begin{tabular}{lcc}
    \toprule
    \textbf{Arquitetura} & \textbf{Cross-Entropy} & \textbf{Focal Loss} \\
    \midrule
    CNN-BiLSTM & E1 & E2 \\
    BiLSTM & E3 & E4 \\
    \bottomrule
  \end{tabular}
  \vspace{.45cm}

  \small
  \begin{itemize}
    \item Perda em cada arquitetura, \textbf{com teste estatístico}: E2 × E1 e E4 × E3.
    \item Etapa convolucional sob a mesma perda, \textbf{descritiva}: E1 × E3 e E2 × E4.
    \item CE sem pesos × FL com $\gamma=2$ e $\alpha_{\mathrm{FA}}=0{,}75$: mede o \textbf{efeito conjunto} da focalização e do peso de classe.
    \item $\alpha_{\mathrm{FA}}=0{,}75$ fica entre a ausência de peso (0,5) e a frequência inversa (0,87); Lin et al. (2017) usaram 0,25 em detecção de objetos, com desbalanceamento extremo.
  \end{itemize}
\end{frame}

\begin{frame}{Protocolo de treinamento}
  \begin{itemize}
    \item \textbf{Otimização:} Adam, taxa de aprendizado $10^{-3}$, lotes de 32 registros, até 50 épocas.
    \item \textbf{Parada antecipada:} 10 épocas sem melhora do F1 de FA na validação.
    \item \textbf{Seleção do modelo:} época com maior F1 de FA na validação; a perda não serve, pois CE e FL têm escalas diferentes.
    \item \textbf{Limiar de decisão:} 0,5 fixo sobre a saída sigmoide; AUC-ROC e AUC-PR independem de limiar.
    \item \textbf{Repetições:} 8 sementes × 4 configurações = \textbf{32 treinamentos}; mesma divisão; teste usado uma única vez.
    \item \textbf{Ambiente:} AMD Ryzen 5 5600; NVIDIA GeForce GTX 1650; 16 GB RAM; SSD de aproximadamente 700 GB.
  \end{itemize}
  {\scriptsize Fonte: Kingma e Ba (2015).}
\end{frame}

\begin{frame}{Avaliação, estatística e custo}
  \small
  \begin{itemize}
    \item \textbf{Desempenho:} F1 de FA como métrica principal; revocação, precisão, especificidade, acurácia, AUC-ROC e AUC-PR como complementares.
    \item \textbf{Hipóteses}, em cada arquitetura: $H_0:\mu_d\leq 0{,}02$ e $H_1:\mu_d>0{,}02$, em que $d_s=F1_{\mathrm{FL},s}-F1_{\mathrm{CE},s}$ é a diferença em cada semente.
    \item \textbf{Teste:} t de Student pareado e unilateral, $\alpha=5\%$. Normalidade por Shapiro--Wilk; se rejeitada, Wilcoxon pareado. Hipótese corroborada só se $H_0$ for rejeitada \textbf{nas duas arquiteturas}.
    \item \textbf{Arquitetura e interação:} descritivas (média, desvio-padrão e IC de 95\%).
    \item \textbf{Custo:} tempo de treinamento, latência por registro e quantidade de parâmetros treináveis.
  \end{itemize}
  {\scriptsize Fontes: Shapiro e Wilk (1965); Wilcoxon (1945).}
\end{frame}

\begin{frame}{Cronograma do TCC2: fevereiro a junho de 2027}
  \centering
  \scriptsize
  \renewcommand{\arraystretch}{1.08}
  \begin{tabular}{@{}p{4.45cm}*{5}{>{\centering\arraybackslash}p{1.4cm}}@{}}
    \toprule
    \textbf{Atividade} & \textbf{Fevereiro} & \textbf{Março} & \textbf{Abril} & \textbf{Maio} & \textbf{Junho} \\
    \midrule
    \textbf{Fase 1: Planejamento} & \barrames & \barrames & & & \\
    Revisão bibliográfica & \barrames & \barrames & & & \\
    Definição do protocolo experimental & \barrames & & & & \\
    \emph{Protocolo definido} & \marcocronograma & & & & \\
    \midrule
    \textbf{Fase 2: Preparação dos dados} & \barrames & \barrames & & & \\
    Obtenção e exploração da base & \barrames & & & & \\
    Pré-processamento dos sinais & \barrames & \barrames & & & \\
    \emph{Base preparada} & & \marcocronograma & & & \\
    \midrule
    \textbf{Fase 3: Implementação dos modelos} & & \barrames & \barrames & & \\
    Implementação da BiLSTM & & \barrames & & & \\
    Implementação da CNN-BiLSTM & & \barrames & \barrames & & \\
    Treinamento das quatro configurações & & & \barrames & & \\
    \emph{Modelos implementados e treinados} & & & \marcocronograma & & \\
    \bottomrule
  \end{tabular}
  \vspace{.16cm}

  \scriptsize Fases 1 a 3 do cronograma do PDF; \textcolor{senacorange}{$\bullet$} indica marco.
\end{frame}

\begin{frame}{Cronograma do TCC2: avaliação e entrega}
  \centering
  \scriptsize
  \renewcommand{\arraystretch}{1.1}
  \begin{tabular}{@{}p{4.45cm}*{5}{>{\centering\arraybackslash}p{1.4cm}}@{}}
    \toprule
    \textbf{Atividade} & \textbf{Fevereiro} & \textbf{Março} & \textbf{Abril} & \textbf{Maio} & \textbf{Junho} \\
    \midrule
    \textbf{Fase 4: Avaliação experimental} & & & \barrames & \barrames & \\
    Avaliação pelas métricas de classificação & & & \barrames & \barrames & \\
    Análise estatística dos resultados & & & & \barrames & \\
    Comparação do custo computacional & & & \barrames & \barrames & \\
    \emph{Experimentos concluídos} & & & & \marcocronograma & \\
    \midrule
    \textbf{Fase 5: Finalização} & & & & \barrames & \barrames \\
    Redação dos resultados e discussão & & & & \barrames & \barrames \\
    Ajustes e retreinamentos, se necessários & & & & \barrames & \barrames \\
    Revisão e formatação do TCC & & & & & \barrames \\
    Preparação da apresentação & & & & & \barrames \\
    \emph{Entrega final} & & & & & \marcocronograma \\
    \bottomrule
  \end{tabular}
  \vspace{.2cm}

  \scriptsize Fases 4 e 5 do cronograma do PDF; \textcolor{senacorange}{$\bullet$} indica marco.
\end{frame}

\section{Considerações finais}

\begin{frame}{Resultados esperados e próximos passos}
  \begin{itemize}
    \item Estimar o efeito de Focal Loss sobre o F1 da classe FA \textbf{em cada arquitetura}.
    \item Descrever se acrescentar a CNN traz ganho diante do tempo e dos parâmetros adicionais.
    \item Executar as \textbf{32 rodadas} do desenho 2 × 2 e avaliar se o ganho médio supera 2 pontos percentuais nas duas arquiteturas.
    \item Reportar as limitações do recorte binário: outros ritmos e sinais ruidosos ficam fora da avaliação.
  \end{itemize}
  \small Os resultados são esperados: a proposta ainda não apresenta modelos treinados nem ganhos comprovados.
\end{frame}

{\setbeamertemplate{footline}{}
\begin{frame}[plain]
  \centering\vfill
  \logosenacgrande\par\vspace{.5cm}
  {\Large\bfseries\color{senacblue}Obrigado!\par}\vspace{.25cm}
  Perguntas?\par\vspace{.55cm}
  {\small Daniel de Araujo Carvalho\\Eduardo Dias Santos}\vfill
\end{frame}}

% Slides de apoio: ficam após o encerramento, para perguntas da banca.
\appendix
\begin{frame}{Apoio: recorte e interpretação}
  \begin{itemize}
    \item CinC 2017: 5.154 registros normais + 771 com FA = 5.925 registros selecionados.
    \item Proporção aproximada: 5.154/771 $\simeq$ 6,7:1.
    \item A exclusão de outros ritmos e registros ruidosos impede comparar o F1 binário diretamente ao escore oficial do desafio.
    \item Métricas calculadas por \textbf{registro}; a classe positiva é FA.
  \end{itemize}
  {\scriptsize Fonte: Clifford et al. (2017).}
\end{frame}

\end{document}
